\documentclass[10pt,a4paper]{amsart}
\usepackage{amsmath,amssymb,amsthm,graphicx}
% Retained minimal preamble for the Stacks Project chapter flat.tex (More on Flatness).
% Verbatim excerpt of the chapter preamble: Xy-pic setup (native preamble lines 14-18) and the native
% theorem declarations with the mathematical macros (native preamble lines 154-217). Copied mechanically;
% nothing added, removed or redefined. The wrapper class and the amsmath/amssymb/amsthm/graphicx base
% packages are supplied by the assembler, so the class/externaldocument part of the native preamble is omitted.
% For commutative diagrams we use Xy-pic
\usepackage[all]{xy}

% We use 2cell for 2-commutative diagrams.
\xyoption{2cell}
% Theorem environments.
%
\theoremstyle{plain}
\newtheorem{theorem}[subsection]{Theorem}
\newtheorem{proposition}[subsection]{Proposition}
\newtheorem{lemma}[subsection]{Lemma}

\theoremstyle{definition}
\newtheorem{definition}[subsection]{Definition}
\newtheorem{example}[subsection]{Example}
\newtheorem{exercise}[subsection]{Exercise}
\newtheorem{situation}[subsection]{Situation}

\theoremstyle{remark}
\newtheorem{remark}[subsection]{Remark}
\newtheorem{remarks}[subsection]{Remarks}

\numberwithin{equation}{subsection}

% Macros
%
\def\lim{\mathop{\mathrm{lim}}\nolimits}
\def\colim{\mathop{\mathrm{colim}}\nolimits}
\def\Spec{\mathop{\mathrm{Spec}}}
\def\Hom{\mathop{\mathrm{Hom}}\nolimits}
\def\Ext{\mathop{\mathrm{Ext}}\nolimits}
\def\SheafHom{\mathop{\mathcal{H}\!\mathit{om}}\nolimits}
\def\SheafExt{\mathop{\mathcal{E}\!\mathit{xt}}\nolimits}
\def\Sch{\mathit{Sch}}
\def\Mor{\mathop{\mathrm{Mor}}\nolimits}
\def\Ob{\mathop{\mathrm{Ob}}\nolimits}
\def\Sh{\mathop{\mathit{Sh}}\nolimits}
\def\NL{\mathop{N\!L}\nolimits}
\def\CH{\mathop{\mathrm{CH}}\nolimits}
\def\proetale{{pro\text{-}\acute{e}tale}}
\def\etale{{\acute{e}tale}}
\def\QCoh{\mathit{QCoh}}
\def\Ker{\mathop{\mathrm{Ker}}}
\def\Im{\mathop{\mathrm{Im}}}
\def\Coker{\mathop{\mathrm{Coker}}}
\def\Coim{\mathop{\mathrm{Coim}}}

% Boxtimes
%
\DeclareMathSymbol{\boxtimes}{\mathbin}{AMSa}{"02}

%
% Macros for moduli stacks/spaces
%
\def\QCohstack{\mathcal{QC}\!\mathit{oh}}
\def\Cohstack{\mathcal{C}\!\mathit{oh}}
\def\Spacesstack{\mathcal{S}\!\mathit{paces}}
\def\Quotfunctor{\mathrm{Quot}}
\def\Hilbfunctor{\mathrm{Hilb}}
\def\Curvesstack{\mathcal{C}\!\mathit{urves}}
\def\Polarizedstack{\mathcal{P}\!\mathit{olarized}}
\def\Complexesstack{\mathcal{C}\!\mathit{omplexes}}
% \Pic is the operator that assigns to X its picard group, usage \Pic(X)
% \Picardstack_{X/B} denotes the Picard stack of X over B
% \Picardfunctor_{X/B} denotes the Picard functor of X over B
\def\Pic{\mathop{\mathrm{Pic}}\nolimits}
\def\Picardstack{\mathcal{P}\!\mathit{ic}}
\def\Picardfunctor{\mathrm{Pic}}
\def\Deformationcategory{\mathcal{D}\!\mathit{ef}}

% --- appended 2026-10-12 (Chief): the project own definitions of the reference/slogan/history
% display environments, copied mechanically from the shared Stacks preamble lines 8-12
% (\usepackage{verbatim} plus the three \newenvironment lines, which the chapter preamble keeps
% verbatim from the project preamble). Nothing added, removed or redefined: these environments
% carry no mathematical content and the project itself typesets them as comments. Without these
% three lines a counted statement whose own text uses them fails with Environment slogan undefined.
\usepackage{verbatim}
\newenvironment{reference}{\comment}{\endcomment}
\newenvironment{slogan}{\comment}{\endcomment}
\newenvironment{history}{\comment}{\endcomment}
\usepackage{cleveref}
\crefname{originallemma}{original lemma}{original lemmas}
\Crefname{originallemma}{Original lemma}{Original lemmas}
\makeatletter
\newlabel{algebra-lemma-spec-functorial}{{00E2}{source}{}{sourcepointer-1}{}}
\newlabel{algebra-lemma-spec-functorial@cref}{{[originallemma][0][]00E2}{[1][0][]source}{}{}{}}
\makeatother
\crefname{originaldefinition}{original definition}{original definitions}
\Crefname{originaldefinition}{Original definition}{Original definitions}
\makeatletter
\newlabel{definition-shrink}{{05H9}{source}{}{sourcepointer-2}{}}
\newlabel{definition-shrink@cref}{{[originaldefinition][0][]05H9}{[1][0][]source}{}{}{}}
\makeatother
\crefname{originallemma}{original lemma}{original lemmas}
\Crefname{originallemma}{Original lemma}{Original lemmas}
\makeatletter
\newlabel{lemma-base-change-one-step-at-x}{{05H8}{source}{}{sourcepointer-3}{}}
\newlabel{lemma-base-change-one-step-at-x@cref}{{[originallemma][0][]05H8}{[1][0][]source}{}{}{}}
\makeatother
\crefname{originallemma}{original lemma}{original lemmas}
\Crefname{originallemma}{Original lemma}{Original lemmas}
\makeatletter
\newlabel{more-morphisms-lemma-finite-morphism-single-point-in-fibre}{{05B8}{source}{}{sourcepointer-4}{}}
\newlabel{more-morphisms-lemma-finite-morphism-single-point-in-fibre@cref}{{[originallemma][0][]05B8}{[1][0][]source}{}{}{}}
\makeatother
\crefname{originallemma}{original lemma}{original lemmas}
\Crefname{originallemma}{Original lemma}{Original lemmas}
\makeatletter
\newlabel{morphisms-lemma-smooth-open}{{056G}{source}{}{sourcepointer-5}{}}
\newlabel{morphisms-lemma-smooth-open@cref}{{[originallemma][0][]056G}{[1][0][]source}{}{}{}}
\makeatother
\title{Shrinking a one step devissage: standard open neighbourhoods of the base, of the middle scheme and of the point}
\author{The Stacks Project Authors}
\date{stacks-project:flat.tex:local-snapshot-20261011; GNU FDL 1.2 or later}
\begin{document}
\maketitle
\noindent\textit{Editorial note.} All delivered mathematical text is the author's own, copied from the pinned Stacks Project chapter file flat.tex; no statement or proof step was rewritten, reordered or removed.\par\medskip
\noindent\textit{Editorial note.} Preamble provenance: the chapter's own retained minimal preamble (sources/stacks-flat/retained-preamble.tex, lines 1-85) is a verbatim excerpt of the native Xy-pic setup and of the native theorem declarations and mathematical macros, followed by the project's own reference/slogan/history display-environment definitions copied mechanically from the shared preamble; the wrapper class and the amsmath/amssymb/amsthm/graphicx base packages are supplied by the assembler, and extra\_packages is left empty because everything needed is already in the retained preamble. No macro was invented or redefined.\par\medskip
\noindent\textit{Editorial note.} Same-chapter pointers are located at their \textbackslash{}label lines inside the pinned chapter file; cross-chapter pointers are located in the project's own tag list staged in the snapshot (sources/stacks-chow/tags.txt), at the verbatim lines carrying their canonical tags.\par\medskip
\noindent\textit{Allowed standard prerequisites (editorial).}\begin{itemize}
\item The notion of a standard shrinking of a one step devissage. Same chapter flat.tex, definition labelled definition-shrink (statement, proof line 972), together with the item labels item-shrink-base, item-shrink-on-Y and item-shrink-on-X declared inside the counted statement. Fixes what must be produced in each of the three statements; the item labels lie inside the counted span, so they need no navigation pointers.
\item Base change of a one step devissage at a point, and functoriality of the spectrum. Same chapter flat.tex, label lemma-base-change-one-step-at-x (proof line 993); Algebra, label algebra-lemma-spec-functorial (proof line 996). Gives part (1): the base changed data still form a one step devissage at the point, and the inverse image of a standard open under a morphism of affine schemes is standard open.
\item Openness of a smooth morphism, and the existence of a standard open neighbourhood of the middle scheme whose preimage lies in a given open. Morphisms, label morphisms-lemma-smooth-open (proof line 1000); More on Morphisms, label more-morphisms-lemma-finite-morphism-single-point-in-fibre (proof line 1016). The first makes the surjectivity argument in part (2) possible, and the second produces the standard open neighbourhood used in part (3) to place the preimage inside the given open U.
\end{itemize}\medskip
\ifdefined\hypertarget\hypertarget{sourcepointer-1}{}\fi
\noindent\textit{Source reference.} Original lemma 00E2: sources/stacks-chow/tags.txt, source lines 519--519. Navigation pointer for \textbackslash{}ref\{algebra-lemma-spec-functorial\} used inside the counted span (line 996); the argument that consumes it is retained in the span.\par\medskip
\ifdefined\hypertarget\hypertarget{sourcepointer-2}{}\fi
\noindent\textit{Source reference.} Original definition 05H9: sources/stacks-flat/pinned-original-source.tex, source lines 954--954. Navigation pointer for \textbackslash{}ref\{definition-shrink\} used inside the counted span (line 972); the argument that consumes it is retained in the span.\par\medskip
\ifdefined\hypertarget\hypertarget{sourcepointer-3}{}\fi
\noindent\textit{Source reference.} Original lemma 05H8: sources/stacks-flat/pinned-original-source.tex, source lines 925--925. Navigation pointer for \textbackslash{}ref\{lemma-base-change-one-step-at-x\} used inside the counted span (line 993); the argument that consumes it is retained in the span.\par\medskip
\ifdefined\hypertarget\hypertarget{sourcepointer-4}{}\fi
\noindent\textit{Source reference.} Original lemma 05B8: sources/stacks-chow/tags.txt, source lines 6876--6876. Navigation pointer for \textbackslash{}ref\{more-morphisms-lemma-finite-morphism-single-point-in-fibre\} used inside the counted span (line 1016); the argument that consumes it is retained in the span.\par\medskip
\ifdefined\hypertarget\hypertarget{sourcepointer-5}{}\fi
\noindent\textit{Source reference.} Original lemma 056G: sources/stacks-chow/tags.txt, source lines 6709--6709. Navigation pointer for \textbackslash{}ref\{morphisms-lemma-smooth-open\} used inside the counted span (line 1000); the argument that consumes it is retained in the span.\par\medskip
\begin{lemma}
\label{lemma-shrink}
With assumption and notation as in
Definition \ref{definition-shrink}
we have:
\begin{enumerate}
\item
\label{item-shrink-base}
If $S' \subset S$ is a standard open neighbourhood of $s$, then
setting $X' = X_{S'}$, $Z' = Z_{S'}$ and $Y' = Y_{S'}$ we obtain a
standard shrinking.
\item
\label{item-shrink-on-Y}
Let $W \subset Y$ be a standard open neighbourhood of $y$.
Then there exists a standard shrinking with $Y' = W \times_S S'$.
\item
\label{item-shrink-on-X}
Let $U \subset X$ be an open neighbourhood of $x$.
Then there exists a standard shrinking with $X' \subset U$.
\end{enumerate}
\end{lemma}

\begin{proof}
Part (1) is immediate from
Lemma \ref{lemma-base-change-one-step-at-x}
and the fact that the inverse image of a standard open under a morphism
of affine schemes is a standard open, see
Algebra, Lemma \ref{algebra-lemma-spec-functorial}.

\medskip\noindent
Let $W \subset Y$ as in (2). Because $Y \to S$ is smooth it is open, see
Morphisms, Lemma \ref{morphisms-lemma-smooth-open}.
Hence we can find a standard open neighbourhood $S'$ of $s$
contained in the image of $W$. Then the fibres of $W_{S'} \to S'$
are nonempty open subschemes of the fibres of $Y \to S$ over $S'$
and hence geometrically irreducible too. Setting $Y' = W_{S'}$
and $Z' = \pi^{-1}(Y')$ we see that $Z' \subset Z$ is a standard open
neighbourhood of $z$. Let $\overline{h} \in \Gamma(Z, \mathcal{O}_Z)$
be a function such that $Z' = D(\overline{h})$. As $i : Z \to X$
is a closed immersion, we can find a function $h \in \Gamma(X, \mathcal{O}_X)$
such that $i^\sharp(h) = \overline{h}$. Take $X' = D(h) \subset X$.
In this way we obtain a standard shrinking as in (2).

\medskip\noindent
Let $U \subset X$ be as in (3). We may after shrinking $U$ assume that
$U$ is a standard open. By
More on Morphisms,
Lemma \ref{more-morphisms-lemma-finite-morphism-single-point-in-fibre}
there exists a standard open $W \subset Y$ neighbourhood of $y$ such
that $\pi^{-1}(W) \subset i^{-1}(U)$. Apply (2) to get a standard
shrinking $X', S', Z', Y'$ with $Y' = W_{S'}$. Since
$Z' \subset \pi^{-1}(W) \subset i^{-1}(U)$ we may replace $X'$ by
$X' \cap U$ (still a standard open as $U$ is also standard open)
without violating any of the conditions defining a standard shrinking.
Hence we win.
\end{proof}

\end{document}
